% =====================================================================
%  Chapter 10: More About Lattice Fermions
%  Gattringer & Lang, "Quantum Chromodynamics on the Lattice" (LNP 788)
%  Reconstructed from extract/ch10.txt; equations auto-numbered.
% =====================================================================

\chapter{More About Lattice Fermions}

The implementation of chiral symmetry on the lattice discussed in Chap.~7
was an unsolved puzzle for many years. The problem has continuously inspired
the lattice community and led to a variety of formulations for fermions on
the lattice. So far we have focused on two kinds of fermions in our
presentation, Wilson-type and overlap fermions. Wilson fermions provide a
simple robust and easy-to-use discretization, while overlap fermions
implement chiral symmetry in the most transparent way.

Here we introduce some of the other ideas for discretizing fermions on the
lattice; most of them rooted in the problem of chiral symmetry. The
exceptions are heavy quark effective formulations designed for simulations
with heavy quarks like charm and bottom. We stress that we can only cover the
basic ideas for the various lattice fermions and refer to the original
literature for more detailed accounts and the current status of numerical
simulations.

\section{Staggered Fermions}

Staggered fermions, often referred to as Kogut--Susskind fermions [1], are a
formulation where the 16-fold degeneracy of the naive discretization is
reduced to only four quarks, while at the same time a remnant chiral symmetry
is maintained. This is achieved by a transformation which mixes spinor and
space--time indices, distributing the quark degrees of freedom on the
hypercubes of the lattice.

\subsection{The Staggered Transformation}

In order to construct staggered fermions we go back to the naive
discretization already presented in (2.29). The free naive fermion action
reads
\begin{equation}
  S_F[\bar\psi, \psi] = a^4 \sum_{n\in\Lambda}
  \left[
    \sum_{\mu=1}^{4} \bar\psi(n)\, \gamma_\mu \,
    \frac{\psi(n+\hat\mu) - \psi(n-\hat\mu)}{2a}
    + m\,\bar\psi(n)\psi(n)
  \right] .
\end{equation}
As we have found in Sect.~5.2, due to the doublers, this action gives rise to
16 quark flavors. In the Wilson Dirac operator 15 of these are removed by
adding the Wilson term at the cost of explicitly breaking chiral symmetry.

However, action (10.1) has a symmetry which can be used to reduce the number
of quark flavors without having to break chiral symmetry. To make explicit
this symmetry one performs a space--time-dependent variable transformation of
the fermion fields $\psi(n)$, $\bar\psi(n)$. This so-called staggered
transformation eliminates the matrices $\gamma_\mu$ in (10.1). To see this we
define new field variables $\psi'(n)$ and $\bar\psi'(n)$ by setting
\begin{equation}
  \psi(n) = \gamma_1^{n_1}\gamma_2^{n_2}\gamma_3^{n_3}\gamma_4^{n_4}
  \,\psi'(n) , \qquad
  \bar\psi(n) = \bar\psi'(n)\,
  \gamma_4^{n_4}\gamma_3^{n_3}\gamma_2^{n_2}\gamma_1^{n_1} .
\end{equation}
The transformation matrices are simply products of the $\gamma_\mu$ raised to
a power given by the corresponding component $n_\mu$ of the site index
$n = (n_1, n_2, n_3, n_4)$. The staggered transformation thus mixes
space--time and Dirac indices.

Since the $\gamma$-matrices obey $\gamma_\mu^2 = 1$, it is obvious that the
mass term is invariant, i.e., $\bar\psi(n)\psi(n) = \bar\psi'(n)\psi'(n)$. In
the kinetic term of (10.1) the $\psi$-field is shifted by one lattice constant
with respect to the $\bar\psi$-field. Thus transformations (10.2) for the
fields on the two sites differ by one power of $\gamma_\mu$. This extra
$\gamma_\mu$ cancels the one which is explicit in the kinetic term of (10.1).
Taking into account the necessary reordering of the $\gamma$-matrices one
finds for the terms of the derivative in, e.g., the 3-direction
\begin{equation}
  \bar\psi(n)\, \gamma_3\, \psi(n \pm \hat{3})
  = (-1)^{n_1 + n_2}\, \bar\psi'(n)\; \mathbf{1} \;
  \psi(n \pm \hat{3})'
\end{equation}
and similar for the terms in the other directions. The action (10.1) turns
into
\begin{equation}
  \begin{split}
    S_F[\bar\psi', \psi'] = a^4 \sum_{n\in\Lambda}
    \Bigl[ &\sum_{\mu=1}^{4} \bar\psi'(n)\; \mathbf{1} \;
      \eta_\mu(n) \frac{\psi'(n+\hat\mu) - \psi'(n-\hat\mu)}{2a}\\
    &\quad + m\,\bar\psi'(n)\psi'(n)
    \Bigr] ,
  \end{split}
\end{equation}
where we have introduced the staggered sign functions
\begin{equation}
  \begin{aligned}
    \eta_1(n) &= 1 , \qquad \eta_2(n) = (-1)^{n_1} , \\
    \eta_3(n) &= (-1)^{n_1 + n_2} , \qquad
    \eta_4(n) = (-1)^{n_1 + n_2 + n_3} ,
  \end{aligned}
\end{equation}
which take over the role of the matrices $\gamma_\mu$. Obviously action (10.4)
is diagonal in Dirac space and has the same form for all four Dirac
components.

The staggered fermion action is obtained by keeping only one of the four
identical components. Coupling the gauge fields we end up with
\begin{multline}
  S_F[\bar\chi, \chi] = a^4 \sum_{n\in\Lambda}
  \Bigg[
    \sum_{\mu=1}^{4} \bar\chi(n)\, \eta_\mu(n)
    \frac{U_\mu(n)\chi(n+\hat\mu) - U_\mu^{\dagger}(n-\hat\mu)\chi(n-\hat\mu)}{2a}
  \\
  \qquad\qquad + m\,\bar\chi(n)\chi(n)
  \Bigg] ,
\end{multline}
where $\chi(n)$ and $\bar\chi(n)$ are Grassmann-valued fields with only color
indices but without Dirac structure. Having discarded 3 of the 4 identical
copies, we can expect that from the 16 quark degrees of freedom of the naive
action only 4 have survived. This hypothesis will be analyzed in the next
section.

We have claimed that the staggered transformation leaves a subset of the
chiral symmetry intact. In order to see this, we first must identify
$\gamma_5$ in the new spinor basis obtained by the transformation (10.2). The
new form of $\gamma_5$ can be read off when transforming a pseudoscalar
bilinear according to (10.2):
\begin{equation}
  \bar\psi(n)\, \gamma_5\, \psi(n) = \eta_5(n)\,
  \bar\psi'(n)\; \mathbf{1} \; \psi'(n) ,
\end{equation}
where we have defined
\begin{equation}
  \eta_5(n) = (-1)^{n_1 + n_2 + n_3 + n_4} .
\end{equation}
This site-dependent sign plays the role of $\gamma_5$ in the staggered world.
It is obvious that for vanishing mass $m$ the staggered action (10.6) is
invariant under the continuous ($\alpha \in \mathbb{R}$) transformation
\begin{equation}
  \chi(n) \longrightarrow e^{i\alpha\eta_5(n)}\chi(n) , \qquad
  \bar\chi(n) \longrightarrow \bar\chi(n)\, e^{i\alpha\eta_5(n)} .
\end{equation}
This is the global chiral symmetry for the staggered theory, where the quark
degrees of freedom are distributed on the hypercube.

\subsection{Tastes of Staggered Fermions}

The staggered transformation mixes lattice and Dirac indices. This raises
conceptual as well as technical questions. An example for the latter is the
problem of constructing hadron interpolators with definite spin and parity.
For details of this construction we refer the reader to the literature (e.g.,
[2]), although some of the ideas can be found in the subsequent discussion.

In order to understand better the conceptual status of staggered fermions, we
here address the question of how many quarks the staggered action (10.6)
describes and which symmetries relate them. For simplicity this study is done
in the free case ($U_\mu(n) \equiv 1$).

For the analysis of staggered fermions we follow the strategy developed in
[3, 4]. The idea is to group together the 16 sites of a hypercube and to
interpret the corresponding degrees of freedom as 4 species of quarks, each of
them with the familiar 4-spinor structure. Let us assume that the $N_\mu$ are
even, i.e., we have an even number of sites in all directions. We consider
the non-intersecting hypercubes with origins separated by multiples of two
lattice spacings. We write our site labels $n_\mu = 0, 1, \ldots, N_\mu - 1$
in terms of labels $h_\mu$ for these hypercubes and labels $s_\mu$ for the
corners of the hypercube as
\begin{equation}
  n_\mu = 2\,h_\mu + s_\mu
  \qquad \text{with} \qquad
  h_\mu = 0, 1, \ldots N_\mu/2 - 1 , \qquad
  s_\mu = 0, 1 .
\end{equation}
With this labeling of sites, the staggered sign function $\eta_\mu(n)$ defined
in (10.5) becomes independent of $h$ and is a function of only the vector $s$:
\begin{equation}
  \eta_\mu(n) = \eta_\mu(2h + s) = \eta_\mu(s) .
\end{equation}
We now want to combine all degrees of freedom, $\chi(2h+s)$, $\bar\chi(2h+s)$
($s_\mu = 0, 1$), sharing a common hypercube label $h$ into a new set of
fields $q(h)_{ab}$, $\bar q(h)_{ab}$ with indices $a, b = 1, 2, 3, 4$. In
order to sneak in the familiar $\gamma$-matrices we define $s$-dependent
$4\times 4$ matrices $\Gamma^{(s)}$ as products of the $\gamma_\mu$ (compare
(10.2)),
\begin{equation}
  \Gamma^{(s)} = \gamma_1^{s_1}\gamma_2^{s_2}\gamma_3^{s_3}\gamma_4^{s_4} .
\end{equation}
It is easy to check that they obey the orthogonality and completeness
relations familiar from Fierz transformation,
\begin{equation}
  \frac{1}{4}\,\tr\left[\Gamma^{(s)\dagger}\Gamma^{(s')}\right]
  = \delta_{s,s'} , \qquad
  \frac{1}{4}\sum_s \Gamma^{(s)*}_{ba}\,\Gamma^{(s)}_{b'a'}
  = \delta_{a,a'}\,\delta_{b,b'} .
\end{equation}
We define new quark fields $q(h)$ and $\bar q(h)$ as linear combinations of
the fields $\chi(2h+s)$, $\bar\chi(2h+s)$:
\begin{equation}
  q(h)_{ab} \equiv \frac{1}{8}\sum_s \Gamma^{(s)}_{ab}\,\chi(2h+s) , \qquad
  \bar q(h)_{ab} \equiv \frac{1}{8}\sum_s
  \bar\chi(2h+s)\,\Gamma^{(s)*}_{ba} .
\end{equation}
Using (10.13), the linear transformations can be inverted to obtain
\begin{equation}
  \chi(2h+s) = 2\,\tr\left[\Gamma^{(s)\dagger}q(h)\right] , \qquad
  \bar\chi(2h+s) = 2\,\tr\left[\bar q(h)\Gamma^{(s)}\right] .
\end{equation}
With the help of these equations we can express the mass term of (10.6) in
terms of the $q$ and $\bar q$ fields:
\begin{equation}
  \begin{split}
    a^4\sum_n \bar\chi(n)\chi(n)
    &= a^4\sum_h\sum_s \bar\chi(2h+s)\,\chi(2h+s) \\
    &= 4a^4\sum_h\sum_s
       \bar q(h)_{ba}\,\Gamma^{(s)}_{ab}\,\Gamma^{(s)*}_{b'a'}\,q(h)_{a'b'} \\
    &= (2a)^4\sum_h \tr\left[\bar q(h)\,q(h)\right] ,
  \end{split}
\end{equation}
where in the second step the completeness relation of (10.13) was used. The
kinetic term of (10.6) is somewhat more tricky since the shifted fields
$\chi(2h + s \pm \hat\mu)$ mix contributions from different hypercubes. We
first must assign the proper hypercube to the individual fields
$\chi(2h + s \pm \hat\mu)$ and only after that we can represent them in terms
of the $q(h)_{ab}$. For the forward shifted field one finds
\begin{equation}
  \begin{aligned}
    \chi(2h + s + \hat\mu)
    &= 2\,\tr\left[\Gamma^{(s+\hat\mu)\dagger}q(h)\right]
    && \text{for } s_\mu = 0 , \\
    \chi(2h + s + \hat\mu)
    &= \chi(2(h+\hat\mu) + s-\hat\mu)
    = 2\,\tr\left[\Gamma^{(s-\hat\mu)\dagger}q(h+\hat\mu)\right]
    && \text{for } s_\mu = 1 .
  \end{aligned}
\end{equation}
From definition (10.12) of the $\Gamma^{(s)}$ follows
$\Gamma^{(s\pm\hat\mu)} = \eta_\mu(s)\,\gamma_\mu\,\Gamma^{(s)}$, which brings
(10.17) into the form
\begin{equation}
  \chi(2h + s + \hat\mu) =
  2\eta_\mu(s)\,\tr\left[\Gamma^{(s)\dagger}\gamma_\mu
  \left(q(h)\delta_{s_\mu,0} + q(h+\hat\mu)\delta_{s_\mu,1}\right)\right] .
\end{equation}
Performing the equivalent steps for the field shifted in negative
$\mu$-direction, $\chi(2h + s - \hat\mu)$, we can write the $\mu$-contribution
to the kinetic term of (10.6) as
\begin{multline}
  4a^4\sum_h\sum_s \frac{1}{2a}
  \tr\left[\bar q(h)\Gamma^{(s)}\right] \\
  \times \tr\left[\Gamma^{(s)\dagger}\gamma_\mu
  \left(q(h)\delta_{s_\mu,0} + q(h+\hat\mu)\delta_{s_\mu,1}
  - q(h-\hat\mu)\delta_{s_\mu,0} - q(h)\delta_{s_\mu,1}\right)\right] .
\end{multline}
Unfortunately we cannot yet sum over $s$ and employ the completeness relation
of (10.13), since in the second factor $\Gamma^{(s)\dagger}$ is combined with
$s$-dependent terms. This problem is overcome by realizing that in the
summation over hypercubes, defined in (10.10), we could have started also at
$\hat\mu$ instead of the origin. If one writes the sums over $h$ and $s$ in
(10.16) with this shifted convention, some of the $q$ and $\bar q$ are shifted
by $\pm\hat\mu$ and the hypercube indices $s_\mu = 0$ and $s_\mu = 1$ are
interchanged. The latter change provides the $s_\mu$ contributions missing in
(10.19). Thus we can average the two ways of writing (10.19) and obtain after
some algebra
\begin{equation}
  \begin{split}
    S_F[\bar q, q] = (2a)^4\sum_h \Big[
    & m\,\tr\left[\bar q(h)q(h)\right] \\
    &+ \sum_\mu \tr\left[\bar q(h)\gamma_\mu\nabla_\mu q(h)\right]
    - a\sum_\mu \tr\left[\bar q(h)\gamma_5\Delta_\mu q(h)\gamma_\mu\gamma_5\right]
    \Big] ,
  \end{split}
\end{equation}
where we have defined derivative operators on the blocked lattice (note that
here the sites are separated by $b \equiv 2a$),
\begin{equation}
  \nabla_\mu f(h) = \frac{f(h+\hat\mu) - f(h-\hat\mu)}{2b} , \qquad
  \Delta_\mu f(h) = \frac{f(h+\hat\mu) - 2f(h) + f(h-\hat\mu)}{b^2} .
\end{equation}
Guided by the form of the first two terms of (10.20), we identify Dirac
indices $\alpha$ and quark species labels $t$ for our fermion fields by
setting
\begin{equation}
  \psi^{(t)}(h)_\alpha \equiv q(h)_{\alpha t} , \qquad
  \bar\psi^{(t)}(h)_\alpha \equiv \bar q(h)_{t\alpha} .
\end{equation}
Expressed in terms of these new fields the free staggered action reads
\begin{multline}
  S_F[\bar\psi, \psi] = b^4 \sum_h \Bigg[
  \sum_{t=1}^{4} m\,\bar\psi^{(t)}(h)\psi^{(t)}(h)
  + \sum_{t=1}^{4}\sum_{\mu=1}^{4}
  \bar\psi^{(t)}(h)\gamma_\mu\nabla_\mu\psi^{(t)}(h) \\
  - \frac{b}{2} \sum_{t,t'=1}^{4}\sum_{\mu=1}^{4}
  \bar\psi^{(t)}(h)\gamma_5(\tau_5\tau_\mu)_{tt'}\,
  \Delta_\mu\psi^{(t')}(h) \Bigg] .
\end{multline}
Here we have introduced the matrices $\tau_\mu = \gamma_\mu^{T}$, $\mu = 1, 2,
\ldots, 5$, and $b$ is the lattice spacing on the lattice of hypercubes. In
order to distinguish the $N_t = 4$ different species of quarks labeled by
$t = 1, 2, 3, 4$ from usual flavor, they are referred to as tastes of
staggered fermions. The first two terms in (10.23) are diagonal in taste space
and represent the mass and kinetic terms for the four tastes of fermions
expected to be described by (10.6). The third term is reminiscent of a Wilson
term, but mixes the different tastes. This taste symmetry-breaking term
reduces the symmetry of the kinetic term which is invariant under independent
vector and axial rotations for each of the four tastes. The taste-breaking
term is only invariant under the remaining symmetry $U(1) \times U(1)$ given
by the rotations
\begin{equation}
  \begin{split}
    \psi' &= e^{i\alpha}\psi , \qquad \psi' = e^{i\beta\Gamma_5}\psi , \\
    \bar\psi' &= \bar\psi\, e^{-i\alpha} , \qquad
    \bar\psi' = \bar\psi\, e^{i\beta\Gamma_5} ,
  \end{split}
\end{equation}
where we have defined the taste-mixing generator $\Gamma_5 = \gamma_5 \otimes
\tau_5$. This symmetry may be identified with a subgroup of the axial taste
symmetry group $SU(N_t)_A$.

\subsection{Developments and Open Questions}

The symmetry of the four tastes, described by the staggered action, is not the
full symmetry of four independent (massless) flavors of quarks, but instead
reduced to (10.24). Although the taste-breaking term vanishes in a naive
continuum limit, its presence in numerical simulations which are necessarily
done at finite lattice spacing raises several questions.

On a more fundamental level one has to understand whether effects of taste
breaking survive the continuum limit as it is performed in an actual
calculation. As we have discussed in Sect.~3.5.4, the ``true continuum
limit'' is obtained by driving the couplings to a critical point where the
correlation length diverges. The open question thus is whether the
taste-breaking terms change the behavior of the theory at the critical point,
or in more fancy words, whether the taste-breaking terms change the
universality class of the theory.

Ignoring these more fundamental issues one can try to minimize the effect of
the taste-breaking terms. In order to come up with a strategy to do so, it is
important to remember that the different tastes are combinations of the
staggered field variables $\chi$ and $\bar\chi$ on a hypercube. Thus the
different tastes see different link variables $U_\mu(n)$. Consequently
strongly fluctuating gauge links will give rise to large taste
symmetry-breaking effects. A possible strategy to ameliorate these effects is
to use gauge actions that suppress such strong fluctuations, see, e.g.,
[5--7]. An alternative approach, already addressed in Sect.~6.2.6, is to apply
blocking or smearing to the gauge fields which locally averages out large
fluctuations [8, 9].

The effect of such blocking or smearing steps can, e.g., be assessed by
analyzing the spectrum of the Dirac operator for staggered fermions (given
here for the $\chi$-basis (10.6)):
\begin{equation}
  D_{st}(n|m) = m\,\delta_{n,m} +
  \sum_{\mu=1}^{4} \eta_\mu(n)
  \frac{U_\mu(n)\delta_{n+\hat\mu,m}
  - U_\mu(n-\hat\mu)^{\dagger}\delta_{n-\hat\mu,m}}{2a} .
\end{equation}
It is easy to see that for the massless case this Dirac operator is
anti-hermitian and obeys a staggered $\gamma_5$-hermiticity:
\begin{equation}
  D_{st}(n|m)^{\dagger} = -D_{st}(n|m) =
  \eta_5(n)\,D_{st}(n|m)\,\eta_5(m) .
\end{equation}
These two properties guarantee that the eigenvalues come in complex conjugate
pairs and all have the same real part given by the quark mass. Thus
fluctuations of the real part of the eigenvalues are impossible and no
exceptional configurations (compare the discussion in Sect.~6.2.5) can occur.
Consequently the determinant (product of all eigenvalues) is always real and
nonnegative and, for nonzero quark mass, strictly positive.

Furthermore, if taste symmetry breaking was absent, then each eigenvalue would
be 4-fold degenerate. In a systematic comparison [10, 11] it was indeed found
that a sufficient amount of blocking or smearing drives the eigenvalues toward
4-fold degeneracy.

Staggered fermions are widely used for dynamical simulations. The reason is
that due to the reduced number of degrees of freedom (no Dirac structure)
staggered fermions are numerically cheaper to simulate and at the same time
are chirally symmetric. However, a problem is that the action describes four
tastes of quarks, while in a realistic QCD simulation one would like to have
two light mass-degenerate $u$ and $d$ quarks and one heavier strange quark. In
order to suitably reduce the number of degrees of freedom it has been proposed
to simulate QCD with an effective action given by
\begin{equation}
  \exp(-S_{eff}) = \exp(-S_G)\,
  \det\left[D_{st}(m_{ud})\right]^{\half}\,
  \det\left[D_{st}(m_s)\right]^{\frac{1}{4}} ,
\end{equation}
where $m_{ud}$ is the average of $u$ and $d$ quark masses and $m_s$ the
strange quark mass. From a mathematical point of view, taking the square or
quartic roots of the determinant is unproblematic, since as we have shown it
is real and positive. However, this procedure is quite nontrivial from a
conceptual perspective and we must ask ourselves if the universality class
remains the same in this approach. Probably even more important is the
question whether the effective action can be expressed in the form of a local
lattice field theory. For a snapshot of the ongoing debate about these issues
see [12--16] and references therein. Although the conceptual problems are not
all resolved, simulations with staggered fermions have found good agreement
with experimental results. Examples are found in [17--20].

\section{Domain Wall Fermions}

The overlap formulation of lattice QCD, discussed in Sect.~7.4, is related to
another version of chiral lattice QCD, so-called domain wall fermions. Domain
wall fermions make use of a 5D lattice and construct chiral Dirac fermions on
a 4D interface of the 5D lattice. The action of the 5D theory is simple,
i.e., rather similar to the usual Wilson formulation. This implies that
already established numerical methods can be applied with only minor
adaptions.

\subsection[Formulation of Lattice QCD with Domain Wall Fermions]{Formulation of Lattice QCD with\\ Domain Wall Fermions}

The basic concepts for domain wall fermions were outlined in the seminal paper
[21] by Kaplan. The ideas were developed further in [22--24] giving rise to
the formulation of domain wall fermions mainly used now.

The idea is to work with a 5D lattice $\Lambda_5 = \Lambda \times
\mathbb{Z}_{N_5}$, where $\Lambda$ is our 4D lattice and $N_5$ denotes the
number of lattice points in the auxiliary 5-direction. On this lattice we use
fermion fields which are described by the 4-component spinors
\begin{equation}
  \begin{aligned}
    &\Psi(n, s)^{\alpha}_c , \quad \bar\Psi(n, s)^{\alpha}_c
    \qquad \text{with} \quad n \in \Lambda , \quad
    s = 0, \ldots, N_5 - 1 , \\
    &\alpha = 1, \ldots, 4 , \quad c = 1, 2, 3 .
  \end{aligned}
\end{equation}
Obviously the Dirac (index $\alpha$) and color (index $c$) structures are
unchanged and we have simply added an additional direction which we label with
$s$. The fermion action for the 5D domain wall theory is given by (setting
$a \equiv 1$)
\begin{equation}
  S_F^{dw}[\Psi, \bar\Psi, U] =
  \sum_{n,m\in\Lambda}\sum_{s,r=0}^{N_5-1}
  \bar\Psi(n, s)\, D^{dw}(n, s|m, r)\, \Psi(m, r) ,
\end{equation}
where we use vector/matrix notation for the Dirac and color indices. The 5D
domain wall Dirac operator $D^{dw}$ is given by
\begin{equation}
  D^{dw}(n, s|m, r) = \delta_{s,r}\, D(n|m) + \delta_{n,m}\, D_5^{dw}(s|r) .
\end{equation}
The first term is diagonal in the auxiliary direction and is built with the
usual 4D Wilson Dirac operator (compare (5.51) where we write explicitly also
color and Dirac indices)
\begin{equation}
  D(n|m) = (4 - M_5)\,\delta_{n,m} -
  \frac{1}{2}\sum_{\mu=\pm1}^{\pm4}
  (1 - \gamma_\mu)\, U_\mu(n)\, \delta_{n+\hat\mu,m} .
\end{equation}
We have introduced a new mass parameter $M_5$, which is not to be confused
with the quark mass parameter $m$ of the 4D theory. The avoidance of doublers
and the positivity of the transfer matrix restrict this parameter to
$0 < M_5 < 1$. The link variables $U_\mu(n)$, $\mu = 1, 2, 3, 4$ are elements
of the gauge group as usual. They do not depend on the coordinate in the fifth
dimension, i.e., identical copies are used for the different 4D slices.

The second contribution to (10.30) is diagonal in $\Lambda$ and the operator
$D_5^{dw}(s|r)$, acting in the 5-direction, is given by
\begin{equation}
  \begin{split}
    D_5^{dw}(s|r) = \delta_{s,r} &- (1-\delta_{s,N_5-1})\, P_-\, \delta_{s+1,r}
    - (1-\delta_{s,0})\, P_+\, \delta_{s-1,r} \\
    &+ m\left(P_-\,\delta_{s,N_5-1}\,\delta_{0,r}
    + P_+\,\delta_{s,0}\,\delta_{N_5-1,r}\right) .
  \end{split}
\end{equation}
Here $P_\pm = (1\pm\gamma_5)/2$ are the usual chiral projectors acting on the
Dirac indices. The parameter $m$ will turn out to be the mass parameter of the
4D target theory. For the fermions at mass $m = 0$ the boundary conditions in
the 5-direction are fixed, since hopping from $s = N_5 - 1$ in positive
5-direction is blocked by the factor $(1-\delta_{s,N_5-1})$, and the hopping
from $s = 0$ in negative direction is eliminated by $(1-\delta_{s,0})$. Only
the terms containing the mass $m$ connect the slices at $s = 0$ and
$s = N_5 - 1$. On the 4D lattice $\Lambda$ the conventional boundary
conditions are used: periodic in the spatial directions 1,2,3 and antiperiodic
in time (4-direction). The link variables $U_\mu(n)$ are periodic in all four
directions of $\Lambda$ and, as mentioned, one uses identical copies for all
values of the 5-coordinate $s$. We stress that in the 5-direction there is no
component $U_5$.

Having presented the action for the 5D fermion fields $\Psi, \bar\Psi$, we can
now construct the 4D physical fields $\psi, \bar\psi$ which live on the 4D
boundary of $\Lambda_5$. These are defined as
\begin{equation}
  \psi(n) = P_-\,\Psi(n, 0) + P_+\,\Psi(n, N_5 - 1) , \qquad
  \bar\psi(n) = \bar\Psi(n, N_5 - 1)\, P_- + \bar\Psi(n, 0)\, P_+ ,
\end{equation}
where $n \in \Lambda$. The physical fields $\psi$ and $\bar\psi$ thus are built
from the degrees of freedom on the first ($s = 0$) and last
($s = N_5 - 1$) 4D slice. The spinors $\psi, \bar\psi$ inherit their Dirac and
color indices from $\Psi, \bar\Psi$.

The physical fields $\psi$ and $\bar\psi$ can now be used to construct the
physical observables of interest. For example the 4D scalar density
$\bar\psi(n)\psi(n)$ assumes the form
\begin{equation}
  \bar\psi(n)\psi(n) = \bar\Psi(n, N_5 - 1)\, P_-\, \Psi(n, 0)
  + \bar\Psi(n, 0)\, P_+\, \Psi(n, N_5 - 1) ,
\end{equation}
when (10.33) is used to express the physical density $\bar\psi(n)\psi(n)$ in
terms of the 5D fields. We remark that (10.34) corresponds exactly to the term
in (10.32) which is proportional to $m$. Thus $m$ can be identified as the
mass parameter of the 4D theory.

The transition to the 4D world can be implemented in a transparent way by
using generating functionals (compare (5.32)). For our purpose we define the
generating functional as
\begin{multline}
  W[J, \bar J] =
  \bigg\langle \exp\Big[\sum_{n\in\Lambda}
  \big(\bar J(n)\psi(n) + \bar\psi(n)J(n)\big)\Big]\bigg\rangle \\
  = \bigg\langle \exp\Big[\sum_{n\in\Lambda}
  \bar J(n)\big(P_-\,\Psi(n, 0) + P_+\,\Psi(n, N_5 - 1)\big) \\
  \qquad
  + \big(\bar\Psi(n, N_5 - 1)\, P_- + \bar\Psi(n, 0)\, P_+\big) J(n)
  \Big]\bigg\rangle_5 .
\end{multline}
Here $J(n)$ and $\bar J(n)$ are Grassmann-valued source fields which have the
same indices (Dirac and color) as the physical fields $\psi(n)$, $\bar\psi(n)$
they couple to. In this equation $\langle\ldots\rangle$ denotes the
expectation value of the 4D target theory and $\langle\ldots\rangle_5$ denotes
the vacuum expectation value in the 5D world, which we discuss in detail in
the next section.

Once the generating functional $W[J, \bar J]$ is defined, one can work with it
like in the regular 4D formulation. In particular $W[J, \bar J]$ can be
differentiated with respect to individual components of $J(n)$ or $\bar J(n)$
which bring the corresponding components of the physical fields $\psi,
\bar\psi$ down from the exponent. After the physical observables are built
from these, one sets the sources to $J = \bar J = 0$ (compare Sect.~5.1.6). As
a result the vacuum expectation values of the 4D target theory are formulated
as the expectation values in the 5D formulation. The latter is then used to
evaluate the expressions.

\subsection{The 5D Theory and Its Equivalence to 4D Chiral Fermions}

So far we have introduced the fields $\Psi$ and $\bar\Psi$ which live on our
5D lattice and their action and discussed how to construct from them the
fields of the 4D theory. We still need to present the detailed definition of
the 5D vacuum expectation values $\langle\ldots\rangle_5$ and finally we
should convince ourselves that for $m = 0$ the 5D theory gives rise to 4D
chiral fermions.

In order to make connection to a 4D overlap-type theory, in addition to the 5D
fermion fields $\Psi, \bar\Psi$ the 5D path integral also contains the fields
$\Phi, \bar\Phi$ which have the same indices as the fermion fields (10.28) but
are bosonic variables (not Grassmann numbers). Thus they are often referred to
as pseudofermion fields (see Sect.~8.1.3) or Pauli--Villars fields. From a
physical point of view their role is the removal of heavy degrees of freedom
which appear for large $N_5$.

The definition of the 5D path integral for the vacuum expectation value
$\langle\ldots\rangle_5$, needed for the generating functional $W[J, \bar J]$
given in (10.35), reads
\begin{equation}
  \langle O\rangle_5 = \frac{1}{Z}
  \int D[\Psi, \bar\Psi, \Phi, \bar\Phi, U]\;
  e^{-S_F^{dw}[\Psi,\bar\Psi,U] - S_B^{pf}[\Phi,\bar\Phi,U] - S_G[U]}\;
  O[\Psi, \bar\Psi, U] .
\end{equation}
$S_F^{dw}[\Psi, \bar\Psi, U]$ is the domain wall action and $S_G[U]$ is some
lattice version of the 4D gauge action, e.g., the Wilson action (2.49). The
path integral is over the 5D fields $\Psi$ and $\bar\Psi$ and the 4D gauge
fields $U$. We have also introduced the integral over the bosonic
pseudofermions $\Phi$ and $\bar\Phi$. The action for the pseudofermions is
given by
\begin{equation}
  \begin{split}
    S_F^{pf}[\Phi, \bar\Phi, U] &=
    \sum_{n,m\in\Lambda}\sum_{s,r=0}^{N_5-1}
    \bar\Phi(n, s)\, D^{pf}(n, s|m, r)\, \Phi(m, r) , \\
    \text{with} \qquad D^{pf}(n, s|m, r) &= \delta_{s,r}\, D(n|m)
    + \delta_{n,m}\, D_5^{pf}(s|r) .
  \end{split}
\end{equation}
The first part is identical to the term in the Dirac operator (10.30), i.e.,
given by the Wilson operator (10.31). Concerning the difference operator
$D_5^{pf}$ for the 5-direction different choices are possible. A variant
convenient for numerical simulations is [25]
\begin{equation}
  D_5^{pf}(s|r) = \left. D_5^{dw}(s|r) \right|_{m=1} .
\end{equation}
A direct connection of the 5D theory (10.36) to a variant of the overlap
formulation was given in [26]. Using an elegant chain of transformations
Neuberger showed (for a slightly different choice of $D^{pf}$) that (see also
[27])
\begin{equation}
  \det D^{dw} = \det D_{N_5}^{ov}\, \det D^{pf} .
\end{equation}
Applying this result to our path integral (10.36) we find
\begin{equation}
  \int D[\Psi, \bar\Psi, \Phi, \bar\Phi]\;
  e^{-S_F^{dw}[\Psi,\bar\Psi,U] - S_B^{pf}[\Phi,\bar\Phi,U]}
  = \frac{\det[D^{dw}]}{\det[D^{pf}]}
  = \det\left[D_{N_5}^{ov}\right] .
\end{equation}
Here we have used that for the path integral over the bosonic pseudofermions
the determinant appears in the denominator (compare Sect.~5.1.4). Obviously
the pseudofermion fields in (10.36) are needed to cancel the factor
$\det[D^{pf}]$ in (10.39). This step can also be viewed as a transformation of
the integration variables with $\det[D^{pf}]$ being the corresponding
Jacobian.

The Dirac operator $D_{N_5}^{ov}$ is often referred to as the truncated
overlap operator and is given by (we drop an irrelevant overall factor of 1/2
compared to the notation in [26])
\begin{equation}
  D_{N_5}^{ov} = 1 + \gamma_5\, \tanh\left(N_5\, \tilde H\right)
  \xrightarrow{N_5\to\infty}
  1 + \gamma_5\, \operatorname{sign}[H] .
\end{equation}
The operator $\tilde H$ is a (nonlocal) variant of the operator $H$ which we
used in Sect.~7.4 to construct the overlap operator (see [26] for the exact
definition). Equation (10.41) shows that in the limit $N_5 \to \infty$, where
the tanh approaches the sign function, the operator assumes the form of the
overlap operator (7.78), and a quick calculation similar to (7.82) shows that
$D_{\infty}^{ov}$ obeys the Ginsparg--Wilson equation (7.29).

The truncated overlap operator describes the 4D theory and in the limit
$N_5 \to \infty$ gives rise to massless fermions. For finite $N_5$ the
chirality-violating effects are exponentially suppressed with an exponent
$\propto N_5$. In practical simulations typical values used are $N_5 = 10$--$30$.
The parameter $M_5$ can be tuned to minimize the chirality-violating effects.

Working with domain wall fermions has the big advantage that numerical
algorithms for 4D Wilson fermions can be easily adapted for the 5D domain wall
formulation. As the Wilson formulation, the domain wall operator also contains
only nearest neighbor terms and no evaluation of matrix-valued functions as
for the overlap operator (see Sect.~7.4.3) is required. For an example of
dynamical simulations with domain wall fermions, see [28].

\section{Twisted Mass Fermions}

Twisted mass QCD (tmQCD) is a formulation which in its simplest form is for
QCD with two mass-degenerate quark flavors of Wilson fermions (QCD with
isospin). The isospin degree of freedom is used to introduce an additional
mass term with a nontrivial isospin structure. This twisted mass term provides
a useful infrared regulator and furthermore can be utilized to obtain O(a)
improvement of the lattice formulation. Twisted mass QCD was first outlined in
[29--31]. In this section we give a brief introduction. For more detailed
recent reviews, which cover also the present status of the numerical analysis
of tmQCD, we recommend [32, 33].

\subsection{The Basic Formulation of Twisted Mass QCD}

We begin our discussion with presenting the defining equations. As already
stated, tmQCD is a theory for two mass-degenerate quarks.\footnote{For the
generalization of tmQCD to the case of quarks with different masses we refer
the reader to [34, 35].} Thus from now on we denote by $\chi$ and $\bar\chi$
quark fields which in addition to Dirac and color indices carry also a flavor
index which may assume $N_f = 2$ values. For notational convenience we will
suppress the indices and use matrix/vector notation instead. The fermion
action for lattice tmQCD with Wilson fermions then reads
\begin{equation}
  S_F^{tw}[\chi, \bar\chi, U] = a^4 \sum_{k,n\in\Lambda}
  \bar\chi(k)\left[D(k|n)\,\mathbb{1}_2 + m\,\mathbb{1}_2\,\delta_{k,n}
  + i\mu\gamma_5\tau^3\,\delta_{k,n}\right]\chi(n) .
\end{equation}
We display unit matrices $\mathbb{1}_2$ only for flavor space and omit them
for color and Dirac. $D(k|n)$ denotes the massless Wilson Dirac operator for a
single flavor (compare (5.51)):
\begin{equation}
  D(k|n) = \frac{4}{a}\,\delta_{k,n} -
  \frac{1}{2a}\sum_{\mu=\pm1}^{\pm4}
  (1 - \gamma_\mu)\, U_\mu(k)\, \delta_{k+\hat\mu,n} .
\end{equation}
Action (10.42) differs from the usual action for two mass-degenerate flavors
with mass $m$ by adding the term $i\mu\gamma_5\tau^3$ to the Dirac operator.
The real parameter $\mu$ is called the twisted mass. We stress, however, that
the conventional mass term is trivial in color, Dirac, and flavor space, while
the twisted mass term is trivial only in color space, has a $\gamma_5$ in
Dirac space, and the third Pauli matrix $\tau^3 = \mathrm{diag}(1,-1)$ acts in
flavor space.

One of the original [29] motivations for introducing the twisted mass
term\footnote{Another motivation being the simplification of renormalization
properties (see the following discussion and Chap.~11).} was its use as an
infrared regulator which removes exceptional configurations (compare
Sect.~6.2.5). Using just the standard mass these are absent only for
sufficiently large values of $m$. The fact that the twisted mass term is a
save remedy against exceptional configurations follows from the following
chain of identities:
\begin{equation}
  \begin{split}
    \det\left[D\,\mathbb{1}_2 + m\,\mathbb{1}_2 + i\mu\gamma_5\tau^3\right]
    &= \det[D + m + i\mu\gamma_5]\; \det[D + m - i\mu\gamma_5] \\
    &= \det[D + m + i\mu\gamma_5]\;
       \det\left[\gamma_5 (D + m - i\mu\gamma_5)\gamma_5\right] \\
    &= \det[D + m + i\mu\gamma_5]\;
       \det[D^{\dagger} + m - i\mu\gamma_5] \\
    &= \det\left[(D + m + i\mu\gamma_5)(D^{\dagger} + m - i\mu\gamma_5)\right] \\
    &= \det\left[(D + m)(D + m)^{\dagger} + \mu^2\right] > 0
       \quad \text{for } \mu \neq 0 .
  \end{split}
\end{equation}
The twisted mass Dirac operator is diagonal in flavor space and thus the
determinant for the 2-flavor operator was written as a product of two
determinants for a single flavor. Subsequently the $\gamma_5$-hermiticity
(5.76) of the Wilson Dirac operator was employed. The inequality in the last
line holds because the eigenvalues of the product $(D + m)(D + m)^{\dagger}$
are real and nonnegative. Thus a nonvanishing value of the twisted mass $\mu$
ensures that the determinant of the twisted mass Dirac operator is strictly
positive and consequently zero eigenvalues, which cause the exceptional
configurations, are excluded. In more detail one can show that the spectrum of
the twisted mass Dirac operator in the complex plane is expelled from a strip
of width $2\mu$ along the real axis [36]. Equation (10.44) establishes that
the determinant is real and positive for arbitrary gauge configurations and
thus standard Monte Carlo techniques are applicable.

The above discussion shows that the twisted mass term may be used as an
alternative infrared regulator which can be combined with the usual mass term.
Since it is possible to work with both, nonvanishing $m$ and $\mu$, it is
convenient to introduce the so-called polar mass $M$ and the twist angle
$\alpha$:
\begin{equation}
  M = \sqrt{m^2 + \mu^2} , \qquad \alpha = \arctan(\mu/m) .
\end{equation}
With these definitions the mass terms of tmQCD can be written as
\begin{equation}
  m\,\mathbb{1}_2 + i\mu\gamma_5\tau^3 = M\, e^{i\alpha\gamma_5\tau^3}
  \qquad \text{with} \qquad
  m = M\cos(\alpha) , \quad \mu = M\sin(\alpha) .
\end{equation}
The case of $\alpha = \pi/2$, i.e., $m = 0$, $\mu > 0$, is often referred to
as maximal or full twist. Below it will be shown that choosing full twist is
special since it implies O(a) improvement. The value $\alpha = 0$ corresponds
to zero twist.

It is instructive to perform a simple transformation of the fermion fields to
new variables $\psi$, $\bar\psi$ defined as
\begin{equation}
  \psi = R(\alpha)\,\chi , \qquad \bar\psi = \bar\chi\, R(\alpha) , \qquad
  R(\alpha) = e^{i\alpha\gamma_5\tau^3/2} .
\end{equation}
A few lines of algebra show that the fermion action (10.42) turns into
\begin{equation}
  S_F[\bar\psi, \psi, U] = a^4 \sum_{k,n\in\Lambda}
  \bar\psi(k)\left[D^{tw}(k|n) + M\,\mathbb{1}_2\,\delta_{k,n}\right]\psi(n) .
\end{equation}
Obviously the twisted mass term has disappeared and is replaced by a
conventional mass term with a mass parameter given by the polar mass $M$
defined in (10.45). The twisted Dirac operator $D^{tw}$ is now a genuine
2-flavor operator and reads
\begin{equation}
  D^{tw}(k|n) = \frac{4\,e^{-i\alpha\gamma_5\tau^3}}{a}\,\delta_{k,n}
  - \frac{1}{2a}\sum_{\mu=\pm1}^{\pm4}
  \left(e^{-i\alpha\gamma_5\tau^3} - \gamma_\mu\right)
  U_\mu(k)\, \delta_{k+\hat\mu,n} .
\end{equation}
We observe that the naive parts of the lattice Dirac operator (the terms with
$\gamma_\mu$) are not affected by the twist and only the Wilson term is
rotated. Thus in the basis $\psi, \bar\psi$, which is referred to as the
physical basis, the mass term and the kinetic terms assume their conventional
form and only the term needed to remove the doublers is introduced in a
twisted form. In the naive continuum limit this latter term is of O(a) and
thus vanishes as $a \to 0$. Our old basis $\chi, \bar\chi$, where the twist
affects the mass term, is known as the twisted basis.

We finally remark that the symmetries of QCD assume a different form in the
twisted basis. To give an example, it is easy to see that the modified parity
transformations (the gauge field still transforms as stated in (5.72))
\begin{equation}
  \chi(n, n_4) \xrightarrow{P} \gamma_4\tau^{1,2}\,\chi(-n, n_4) , \qquad
  \bar\chi(n, n_4) \xrightarrow{P} \bar\chi(-n, n_4)\,\gamma_4\tau^{1,2}
\end{equation}
leave the action (10.42) invariant. Note that we have two different choices
and can use either of the two Pauli matrices $\tau^1$ or $\tau^2$. A collection
of the twisted forms of other symmetries can, e.g., be found in the appendix
of [33].

\subsection{The Relation Between Twisted and Conventional QCD}

Having introduced a new type of infrared regulator, one of course has to
establish that in some limit the twisted formulation describes conventional
QCD. At the end of the last section we have seen that when changing to the
physical basis $\psi, \bar\psi$, the twist affects only the Wilson term which
in the naive continuum limit vanishes anyway. This is already a good sign, but
the equivalence of conventional and tmQCD has to be established also for the
proper continuum limit (compare Sect.~3.5.4).

We begin with studying the relation between conventional and twisted mass QCD
by analyzing the situation in the continuum. The conventional action for two
mass-degenerate fermions in the continuum is given by
\begin{equation}
  S_F[\bar\psi, \psi, A] = \int d^4x\;
  \bar\psi(x)\left(\gamma_\mu D_\mu(x)\,\mathbb{1}_2
  + M\,\mathbb{1}_2\right)\psi(x) ,
\end{equation}
where $D_\mu(x) = \partial_\mu + iA_\mu(x)$. Vacuum expectation values of
some operator $O$ are computed with the path integral which is formally
defined as ($S_G[A]$ denotes the action for the gauge fields)
\begin{equation}
  \langle O\rangle = \frac{1}{Z}
  \int D[\bar\psi, \psi, A]\;
  e^{-S_F[\bar\psi,\psi,A] - S_G[A]}\; O[\bar\psi, \psi, A] .
\end{equation}
Relating this conventional form of the vacuum expectation value to the
corresponding expression in tmQCD is now merely an exercise in changing
integration variables. The transformation we need is the one given in (10.47).
Under this transformation the conventional continuum action (10.51) changes
to the continuum tmQCD action given by
\begin{equation}
  S_F^{tw}[\bar\chi, \chi, A] = \int d^4x\;
  \bar\chi(x)\left(\gamma_\mu D_\mu(x)\,\mathbb{1}_2
  + m\,\mathbb{1}_2 + i\mu\gamma_5\tau^3\right)\chi(x) ,
\end{equation}
where the mass parameters $m$ and $\mu$ are related to the mass parameter $M$
of the conventional action (10.51) through (10.46). Since transformation
(10.47) is non-anomalous, the integration measure remains invariant and,
somewhat formally, we can write $D[\bar\psi, \psi] = D[\bar\chi, \chi]$. Thus
we define vacuum expectation values in the twisted basis as
\begin{equation}
  \langle O\rangle^{tw} = \frac{1}{Z^{tw}}
  \int D[\bar\chi, \chi, A]\;
  e^{-S_F^{tw}[\bar\chi,\chi,A] - S_G[A]}\;
  O^{tw}[\bar\chi, \chi, A] .
\end{equation}
We are left with the task of relating an operator $O[\bar\psi, \psi, A]$ in
the physical basis to its counterpart $O^{tw}[\bar\chi, \chi, A]$ in the
chiral basis. This is a simple exercise as we illustrate in an example, where
we choose $O$ to be the nonsinglet axial--pseudoscalar correlator,
\begin{equation}
  O[\bar\psi, \psi] = A^1_\mu(x)\,P^1(y)
  \qquad \text{where} \qquad
  A^a_\mu = \bar\psi\,\gamma_\mu\gamma_5\,\frac{\tau^a}{2}\,\psi , \qquad
  P^a = \bar\psi\,\gamma_5\,\frac{\tau^a}{2}\,\psi .
\end{equation}
A brief calculation shows that transformation (10.47) gives rise to
\begin{equation}
  \begin{aligned}
    A^{1}_\mu &\to \cos(\alpha)\,A^{1,tw}_\mu - \sin(\alpha)\,V^{2,tw}_\mu ,
    \qquad P^{1} \to P^{1,tw} \qquad \text{with} \\
    A^{a,tw}_\mu &= \bar\chi\,\gamma_\mu\gamma_5\,\frac{\tau^a}{2}\,\chi ,
    \qquad
    V^{a,tw}_\mu = \bar\chi\,\gamma_\mu\,\frac{\tau^a}{2}\,\chi , \qquad
    P^{a,tw} = \bar\chi\,\gamma_5\,\frac{\tau^a}{2}\,\chi .
  \end{aligned}
\end{equation}
Inserting these, we find for our correlator the following relation between the
vacuum expectation value in the conventional and the twisted form:
\begin{equation}
  \begin{split}
    \langle O\rangle
    &= \left\langle A^1_\mu(x)\,P^1(y)\right\rangle \\
    &= \cos(\alpha)\left\langle A^{1,tw}_\mu(x)\,P^{1,tw}(y)\right\rangle
    - \sin(\alpha)\left\langle V^{1,tw}_\mu(x)\,P^{1,tw}(y)\right\rangle \\
    &= \left\langle O^{tw}\right\rangle^{tw} .
  \end{split}
\end{equation}
In an equivalent way arbitrary $n$-point functions of conventional QCD can be
mapped onto linear combinations of $n$-point functions of tmQCD.

The above discussion is based on the (formal) expressions in the continuum. To
carry the arguments for the relation between conventional and tmQCD over to
the lattice one has to consider the renormalized theory in the continuum limit
using a mass-independent renormalization scheme [30]. The renormalized mass
and twisted mass parameters are
\begin{equation}
  m^{(r)} = Z_m\,(m - m_c) , \qquad \mu^{(r)} = Z_\mu\,\mu ,
\end{equation}
and the twist angle in the renormalized theory must be defined as
\begin{equation}
  \alpha = \arctan\left(\mu^{(r)}/m^{(r)}\right) .
\end{equation}
Maximal twist, i.e., $\alpha = \pi/2$, corresponds to $m^{(r)} = 0$. This case
is particularly simple, since setting the bare mass parameter to its critical
value, $m = m_c$, which can, e.g., be identified by a vanishing PCAC mass (see
Sect.~11.1.2), already leads to $m^{(r)} = 0$.

In the renormalized theory we can use the arguments given for the continuum
discussion above and conclude that standard QCD correlation functions can be
expressed as linear combinations of correlators in tmQCD. This relation
remains valid for finite lattice spacing up to discretization errors [30].

\subsection{O(a) Improvement at Maximal Twist}

Although an important initial motivation for the introduction of tmQCD was the
cure for exceptional configurations, today the property of O(a) improvement is
considered more important. We now briefly address this feature of tmQCD at
maximal twist [37].

With the introduction of the twisted mass term we now have two parameters $m$,
$\mu$, which define the physics we want to describe. The relative size of the
two mass parameters is characterized by the twist angle $\alpha$. We have
already observed that among the possible values of $\alpha$, the case of
$\alpha = \pi/2$, i.e., maximal twist is singled out. For maximal twist the
discretization effects of O(a) vanish and the leading corrections appear only
at O(a$^2$). This property is referred to as O(a) improvement of tmQCD and
will be discussed now.

The special role of $\alpha = \pi/2$ can already be seen in the free case. In
momentum space the twisted mass Dirac operator reads (compare (5.48) for the
case of Wilson fermions with conventional mass term)
\begin{equation}
  \frac{i}{a}\sum_{\mu=1}^{4}\gamma_\mu\sin(p_\mu a)
  + \frac{1}{a}\sum_{\mu=1}^{4}(1-\cos(p_\mu a))
  + M\cos(\alpha) + iM\sin(\alpha)\gamma_5\tau^3 ,
\end{equation}
where we have dropped all unit matrices. It is straightforward to check that
the inverse of this matrix, i.e., the propagator in momentum space, is given
by
\begin{equation}
  \frac{-\frac{i}{a}\sum_\mu\gamma_\mu\sin(p_\mu a)
  + \frac{1}{a}\sum_\mu(1-\cos(p_\mu a)) + M\cos(\alpha)
  - iM\sin(\alpha)\gamma_5\tau^3}
  {\left[\frac{1}{a}\sum_\mu\sin(p_\mu a)\right]^2
  + \left[\frac{1}{a}\sum_\mu(1-\cos(p_\mu a)) + M\cos(\alpha)\right]^2
  + M^2\sin(\alpha)^2}
\end{equation}
The energy of the fermion described by this propagator is given by the
position of the pole. Expanding the denominator in $a$ gives
\begin{equation}
  p^2\left(1 + a\,M\cos(\alpha)\right) + M^2 + O(a^2) ,
\end{equation}
and for the two poles we find ($M^2/p^2$ is held fixed)
\begin{equation}
  ip_4 = \pm\sqrt{p^2 + M^2}
  \mp a\cos(\alpha)\,\frac{M^3}{2\sqrt{p^2 + M^2}} + O(a^2) .
\end{equation}
It is evident that the O(a) correction to the dispersion relation in the
continuum, $ip_4 = \pm\sqrt{p^2 + M^2}$, comes with a factor of $\cos(\alpha)$.
This allows to turn off the O(a) term by choosing maximal twist, i.e., by
setting $\alpha = \pi/2$. From an algebraic point of view this result can be
understood by noting that at $\alpha = \pi/2$ the Wilson term and the mass
term, which now consist of only the twisted part, are orthogonal in isospin
space. Thus a mixed term, which is of O(a), cannot emerge.

Having illustrated that the spectrum of the free theory is O(a) improved at
maximal twist is certainly interesting but of course one would like to
establish O(a) improvement for full tmQCD. This problem can be addressed by
performing the Symanzik improvement program (compare Sect.~9.1) for tmQCD. In
order to have maximal twist in the renormalized theory we need a vanishing
renormalized quark mass parameter $m^{(r)}$. Thus we consider the case where
the bare standard mass parameter $m$ is tuned to its critical value,
$m = m_c$, which may be defined by the requirement of a vanishing PCAC quark
mass (compare Sects.~9.1.4 and 11.1.2). Setting $m = m_c$ implies for the
renormalized quark mass parameter, $m^{(r)} = 0$, as needed for maximal twist.
In this case the relevant contribution to the effective continuum action is
(we show only the leading term -- for a complete list see [31])
\begin{equation}
  S_F^{eff} = S_0 + a\,S_1 + \ldots + O(a^2) ,
\end{equation}
with
\begin{equation}
  S_0 = \int d^4x\;
  \bar\chi\left(\gamma_\mu D_\mu + i\mu\gamma_5\tau^3\right)\chi , \qquad
  S_1 = c_{sw}\int d^4x\; \bar\chi\,\sigma_{\mu\nu}F_{\mu\nu}\,\chi .
\end{equation}
The correction $S_1$ to the continuum action $S_0$ is treated as an insertion
in expectation values, such that we obtain for the vacuum expectation value of
some operator $O$:
\begin{equation}
  \langle O\rangle = \langle O\rangle_0 + a\,\langle\Delta O\rangle_0
  - a\,\langle O\,S_1\rangle_0 + O(a^2) .
\end{equation}
The expectation value $\langle\ldots\rangle_0$ is with respect to the action
$S_0$ and by $\Delta O$ we denote the counterterms for the operator $O$. The
operators $O$ may be classified according to their symmetry under the discrete
chirality transformation
\begin{equation}
  \chi \longrightarrow i\gamma_5\tau^1\chi , \qquad
  \bar\chi \longrightarrow i\bar\chi\gamma_5\tau^1 ,
\end{equation}
which is one of the symmetries of massless QCD with two flavors (see
Sect.~7.1.2). Operators can be decomposed into components with definite
transformation properties under (10.67):
\begin{equation}
  O \longrightarrow \pm O \qquad \text{with} \qquad
  \Delta O \longrightarrow \mp\Delta O .
\end{equation}
The terms of the effective action transform as
\begin{equation}
  S_0 \longrightarrow S_0 , \qquad S_1 \longrightarrow -S_1 .
\end{equation}
Using (10.68) and (10.69) we can analyze the contributions to the expectation
values (10.66). For chirally even operators we find the symmetries
\begin{equation}
  \langle O\,S_1\rangle_0 = -\langle O\,S_1\rangle_0 = 0 , \qquad
  \langle\Delta O\rangle_0 = -\langle\Delta O\rangle_0 = 0 ,
\end{equation}
while for odd operators
\begin{equation}
  \langle O\rangle_0 = -\langle O\rangle_0 = 0 , \qquad
  \langle O\,S_1\rangle_0 = \langle O\,S_1\rangle_0 , \qquad
  \langle\Delta O\rangle_0 = \langle\Delta O\rangle_0 .
\end{equation}
Putting things together we conclude
\begin{equation}
  \begin{aligned}
    \langle O\rangle &= \langle O\rangle_0 + O(a^2)
    && \text{for } O \text{ even} , \\
    \langle O\rangle &= a\left(\langle\Delta O\rangle_0
    - \langle O\,S_1\rangle_0\right) + O(a^2)
    && \text{for } O \text{ odd} .
  \end{aligned}
\end{equation}
This implies that either only O(a$^2$) terms survive (for even operators) or
the expectation value $\langle O\rangle$ vanishes in the continuum limit (odd
operators). Thus we can conclude that the vacuum expectation values of
nonvanishing operators are O(a) improved at tmQCD with maximal twist. The only
parameter that has to be tuned to achieve this is the bare quark mass which
has to be set to its critical value, $m = m_c$.

The ease with which O(a) improvement is obtained and the technical advantages
which come with the removal of exceptional configurations make tmQCD an
attractive lattice formulation. From a numerical point of view no new
techniques are required, which, for example, are needed in the case of
dynamical overlap calculations. A rich program of numerical simulations with
tmQCD was conducted in recent years and for a comprehensive recent overview we
refer to [33].

\section{Effective Theories for Heavy Quarks}

For lattice calculations involving heavy quarks, such as the charm and the
bottom quark, special techniques are needed. The central idea is to remove the
dominant scale, the mass $m_h$ of the heavy quark, and to work with an
effective Lagrangian. This procedure gives rise to two formulations known as
nonrelativistic QCD (NRQCD) and heavy quark effective theory (HQET).

\subsection{The Need for an Effective Theory}

We begin with a short discussion of the scales involved when working with
heavy quarks. The masses of the charm and bottom quarks in the
$\overline{\mathrm{MS}}$ scheme [38] are
\begin{equation}
  m_c \approx 1.27(9)\,\mathrm{GeV} , \qquad
  m_b \approx 4.20(12)\,\mathrm{GeV} .
\end{equation}
The masses for typical charmonium or bottomium states are, e.g.,
\begin{equation}
  M_{J/\psi} = 3.096\,\mathrm{GeV} , \qquad
  M_{\Upsilon} = 9.460\,\mathrm{GeV} .
\end{equation}
If we want to reliably describe these states on the lattice the cutoff $1/a$
should be larger than these masses. Using (6.71) for converting MeV to fm, one
finds that the lattice spacing has to obey
\begin{equation}
  a < 0.064\,\mathrm{fm} \quad \text{for charm} , \qquad
  a < 0.021\,\mathrm{fm} \quad \text{for bottom} .
\end{equation}
On the other hand we need a spatial volume that is sufficiently large to
accommodate the hadrons we want to study and a long enough temporal extent to
be able to track the Euclidean propagators when we determine their masses. A
typical value would be a spatial extent of about 2 fm. While for the charm
quark lattices of sufficient size will be in reach soon, for the bottom quark
a naive approach to simulations with heavy quarks would require very large
lattices with typical sizes of at least $\order{100}^4$.

A large part of the mass of the $J/\psi$ and $\Upsilon$ mesons comes from the
valence quark masses themselves. Thus if one is able to remove this trivial
contribution to the hadron mass from the theory, the typical energy scales of
the problem are reduced considerably. Then one can work again with a lattice
spacing which is large enough such that the required number of lattice points
is sufficiently small for a numerical simulation.

\subsection{Lattice Action for Heavy Quarks}

The first ideas for treating heavy quarks on the lattice were put forward in
[39--43]. As for the case of heavy quarks in the continuum (see [44] for a
standard text), the central idea is to scale out the mass $m_h$ of the heavy
quark. This is done in a systematic way by the Foldy--Wouthuysen
transformation well known from quantum mechanics. In the system where the
heavy quark is at rest it provides an expansion of the action for the heavy
flavors in terms of $1/m_h$. Following a standard textbook such as [45] (see
also [42, 44]), the Euclidean continuum action density $L = \bar q\,(\gamma_\mu
D_\mu + m_h)\,q$ with $D_\mu = \partial_\mu + iA_\mu$ is expanded in the form
(4 is the time direction)
\begin{equation}
  \bar q\,(\gamma_\mu D_\mu + m_h)\,q
  \longrightarrow \bar\psi_h\left(m_h + D_4
  - \frac{D^2}{2m_h} - \frac{\boldsymbol{\sigma}\cdot\boldsymbol{B}}{2m_h}
  \right)\psi_h + \order{1/m_h^2} .
\end{equation}
Here, $\boldsymbol{B}$ denotes the chromomagnetic fields and
$D^2 = \sum_{j=1}^{3}D_j^2$ is the spatial covariant Laplace operator. On the
right-hand side the original 4-spinors $q, \bar q$ for the relativistic quarks
are replaced by projected, nonrelativistic spinors $\psi_h, \bar\psi_h$ which
obey
\begin{equation}
  P_+\psi_h = \psi_h , \qquad \bar\psi_h P_+ = \bar\psi_h , \qquad
  P_+ = \half(1 + \gamma_4) .
\end{equation}
It has to be remarked that in expansion (10.76) the usual renormalizability of
QCD is obtained only by including all orders in $1/m_h$. How vacuum
expectation values with only a finite number of terms in (10.76) can be
defined such that they have a finite continuum limit will be described in the
next section.

Let us now invoke the lattice. Guided by expansion (10.76) of the continuum
action we write the lattice action up to a given order $N$ in the expansion as
\begin{equation}
  \begin{split}
    S_{h,N} &= S_{stat} + \sum_{\nu=1}^{N}S^{(\nu)} , \\
    S_{stat} &= \sum_{n\in\Lambda}
    \bar\psi(n)\left(\bar\nabla_4 + \delta m_h\right)\psi(n) , \\
    S^{(\nu)} &= \sum_i \omega_i^{(\nu)} S_i^{(\nu)} .
  \end{split}
\end{equation}
Here we have written explicitly the static contribution $S_{stat}$ which does
not give rise to spatial propagation. In this term the mass $m_h$ was dropped
to get rid of the unwanted scale $m_h$. For dimensional reasons a mass type of
term has to be included, which is taken into account by a residual mass
parameter $\delta m_h$ which is also needed for matching and renormalization
purposes. By $\bar\nabla_4$ we denote the temporal backward derivative on the
lattice, $\bar\nabla_4 f(n) \equiv f(n) - U_4(n-\hat 4)\,f(n-\hat 4)$. The
nonstatic terms are organized in powers of $1/m_h$ labeled by the index
$\nu$. The leading nonstatic terms are (compare (10.76))
\begin{equation}
  \begin{split}
    S_1^{(1)} &= -\frac{1}{2}\sum_{n\in\Lambda}
    \bar\psi_h(n)\,\boldsymbol{\sigma}\cdot\boldsymbol{B}\,\psi_h(n) , \\
    S_2^{(1)} &= -\frac{1}{2}\sum_{n\in\Lambda}
    \bar\psi_h(n)\,D^2\,\psi_h(n) , \\
    S_1^{(2)} &= \frac{1}{8}\sum_{n\in\Lambda}
    \bar\psi_h(n)\,\boldsymbol{\nabla}\cdot\boldsymbol{E}\,\psi_h(n) , \\
    &\qquad \text{etc.}
  \end{split}
\end{equation}
On the lattice, the chromomagnetic field $\boldsymbol{B}$ and the
chromoelectric field $\boldsymbol{E}$ are implemented by a suitable
discretization which expresses these fields through the plaquettes $U_{\mu\nu}$,
e.g.,
\begin{equation}
  E_j = F_{j4} , \qquad B_j = \frac{\epsilon_{j4\mu\nu}}{2}F_{\mu\nu} , \qquad
  F_{\mu\nu} = i\,\mathrm{Im}\,(1 - U_{\mu\nu}) .
\end{equation}
This discretization of the field strength tensor $F_{\mu\nu}$ follows from
expansion (2.53) of the plaquette $U_{\mu\nu}$ used for the construction of
the gauge action. The covariant (spatial) Laplace operator $D^2$ is discretized
as (we set $a \equiv 1$)
\begin{equation}
  D^2 f(n) \longrightarrow
  \sum_{j=1}^{3}\left[U_j(n)f(n+\hat j) - 2f(n)
  + U_j(n-\hat j)^{\dagger}f(n-\hat j)\right] .
\end{equation}
In the classical theory the expansion coefficients $\omega_i^{(\nu)}$ would be
given by $\omega_i^{(\nu)} = 1/m_h^{\nu}$, as can be seen from comparing
(10.76), (10.78), and (10.79). Here we are using HQET as an effective theory
for fully quantized QCD and have to keep the coefficients $\omega_i^{(\nu)}$
as free parameters which will be fixed later when matching HQET and QCD. It
is, however, important to keep in mind the power counting of the coefficients
in expansion (10.78):
\begin{equation}
  \omega_i^{(\nu)} = \order{m_h^{-\nu}}
  \qquad \left(\text{and } \order{a} \equiv \order{1/m}\right) .
\end{equation}
We remark that ultimately we are interested in QCD with both light and heavy
flavors. The action for the light flavors is included in a form discussed in
the previous chapters and will be added in the next section. We stress that
also the observables one is interested in may mix heavy and light quarks. A
typical example would be the study of the heavy--light $D$ and $B$ mesons.

\subsection{General Framework and Expansion Coefficients}

Having discussed the expansion of the action and its discretization, we need
to address the quantization and renormalization of our effective theory. In
other words, the coefficients of the effective action have to be determined.
Several approaches to this problem can be found in the literature. In our
presentation we follow the nonperturbative strategy outlined [46] for HQET and
reviewed in great detail in [47]. For a more general presentation of recent
results in NRQCD and HQET we refer the reader to [48, 49].

We have already pointed out that expanding the action in $1/m_h$ up to a
fixed-order $N$ spoils the usual renormalizability of QCD. However, it can be
shown that the static term $S_{stat}$ of the effective action alone gives rise
to a renormalizable theory. Thus the Boltzmann factor for the heavy quark part
of the action is expanded as
\begin{equation}
  e^{-S_h} = e^{-S_{stat}}
  \left(1 - S^{(1)} - S^{(2)} + \half\left(S^{(1)}\right)^2 + \ldots\right) ,
\end{equation}
and the nonstatic terms appear as insertions in expectation values computed
with the static action. The expansion is terminated after all terms up to the
desired order $N$ are included. In (10.83) all terms up to order $N = 2$ are
displayed.

The form (10.83) of the expansion leads to the formulation usually referred to
as HQET. An alternative is to keep in the exponent also the $D^2$ contribution,
the $S_2^{(1)}$ term in (10.79), together with $S_{stat}$. This leads to the
NRQCD formulation (see, e.g., [50]).

For our final expression we still have to add the action for the light quarks,
$S_L$, and the action for gauge fields, $S_G$. We obtain for the expectation
value of some observable $O$ in HQET on the lattice
\begin{equation}
  \begin{split}
    \langle O\rangle_h &= \frac{1}{Z_h}
    \int D[\Phi]\; e^{-S_G - S_L - S_{stat}}
    \left(1 - S^{(1)} - S^{(2)}
    + \half\left(S^{(1)}\right)^2 + \ldots\right)O , \\
    Z_h &= \int D[\Phi]\; e^{-S_G - S_L - S_{stat}}
    \left(1 - S^{(1)} - S^{(2)}
    + \half\left(S^{(1)}\right)^2 + \ldots\right) .
  \end{split}
\end{equation}
The integration is over all fields collectively denoted by $\Phi$, i.e., the
light and heavy quark fields and the gauge variables. We remark that for many
applications also the observable $O$ is expanded in $1/m_h$. It can be shown
that the expanded static vacuum expectation values in (10.84) can be
renormalized when all local operators with the proper symmetries and dimension
$\leq N$ are included.

The theory defined by (10.84) is described by the following parameters: From
the gauge and light quark sectors we have the gauge coupling, the masses of
the light quarks, and maybe some coefficients of improvement terms. The heavy
quark part contains $\delta m_h$ and the coefficients $\omega_i^{(\nu)}$,
$\nu = 1, \ldots, N$, where $N$ is the desired order of the expansion.

The remaining problem is how to set the parameters of HQET such that
expectation values match their QCD counterparts. In principle it would be
possible to evaluate a set of observables $\Phi_k(m_h)$ in both theories and
to require
\begin{equation}
  \Phi_k^{HQET}(m_h) = \Phi_k^{QCD}(m_h) + \order{m_h^{-(N+1)}} ,
\end{equation}
where the number of observables has to equal the number of parameters.
However, such a brute force procedure faces exactly the difficulties which we
outlined in Sect.~10.4.1, i.e., the need for large, very fine lattices.

A nonperturbative strategy to overcome this obstacle was presented in [46]. The
idea is to start with a sufficiently fine lattice such that the boundaries for
the cutoff given in (10.75) are obeyed. In order to have a lattice which in
lattice units is small enough to be accessible to a numerical simulation, one
is restricted to lattices with a small physical size of
$L = 0.2$--$0.5$ fm. On the small lattice the matching conditions
\begin{equation}
  \Phi_k^{HQET}(m_h, L) = \Phi_k^{QCD}(m_h, L) + \order{m_h^{-(N+1)}}
\end{equation}
are imposed. The HQET part of the calculation is repeated now using a lattice
of extent $2L$. The two results can be used to compute the step-scaling
functions $F_k$ (see [51]) defined as
\begin{equation}
  \Phi_k^{HQET}(m_h, 2L) = F_k\, \Phi_k^{HQET}(m_h, L) .
\end{equation}
The dimensionless functions $F_k$ describe the change of the complete set of
observables under a volume change $L \to 2L$. For the step-scaling functions
the continuum limit can be performed and subsequently they are used to
transport the physical observables to sufficiently large volumes where contact
with experimental results can be made.

We conclude with emphasizing again that, although for the matching we here
mainly follow the proposal [46], several variants of obtaining effective
theories for heavy quarks have been discussed. Various applications of HQET,
e.g., the determination of the mass of the $b$-quark [46, 52], have been
proposed and more general reviews of the corresponding results can be found in
[48--50, 54, 55].

\section*{References}

\begin{enumerate}[label={[\arabic*]}]
  \item J.~Kogut and L.~Susskind: Phys.\ Rev.\ D 11, 395 (1975).
  \item G.~Kilcup and S.~Sharpe: Nucl.\ Phys.\ B 283, 493 (1987).
  \item F.~Gliozzi: Nucl.\ Phys.\ B 204, 419 (1982).
  \item H.~Kluberg-Stern, A.~Morel, O.~Napoly, and B.~Petersson: Nucl.\ Phys.\
        B 220, 447 (1983).
  \item T.~Takaishi: Phys.\ Rev.\ D 54, 1052 (1996).
  \item K.~Orginos, D.~Touissant, and R.~Sugar: Phys.\ Rev.\ D 60, 054503
        (1999).
  \item G.~P.~Lepage: Phys.\ Rev.\ D 59, 074502 (1999).
  \item M.~Albanese et al.: Phys.\ Lett.\ B 192, 163 (1987).
  \item A.~Hasenfratz and F.~Knechtli: Phys.\ Rev.\ D 64, 034504 (2001).
  \item E.~Follana, A.~Hart, and C.~Davies: Phys.\ Rev.\ Lett.\ 93, 241601
        (2004).
  \item S.~D\"urr, C.~Hoelbling, and U.~Wenger: Phys.\ Rev.\ D 70, 094502
        (2004).
  \item S.~D\"urr: PoS LAT2005, 021 (2005).
  \item S.~R.~Sharpe: Phys.\ Rev.\ D 74, 014512 (2006).
  \item M.~Creutz: Phys.\ Lett.\ B 649, 230 (2007).
  \item A.~Kronfeld: PoS LATTICE2007, 016 (2007).
  \item M.~F.~L.~Golterman (unpublished) arXiv:0812.3110 [hep-lat] (2008).
  \item C.~T.~H.~Davies et al.: Phys.\ Rev.\ Lett.\ 92, 022001 (2004).
  \item C.~Aubin et al.: Phys.\ Rev.\ D 70, 114501 (2004).
  \item C.~W.~Bernard et al.: Phys.\ Rev.\ D 64, 054506 (2001).
  \item A.~Bazavov et al. (unpublished) arXiv:0903.3598 [hep-lat] (2009).
  \item D.~B.~Kaplan: Phys.\ Lett.\ B 288, 342 (1992).
  \item Y.~Shamir: Nucl.\ Phys.\ B 406, 90 (1993).
  \item Y.~Shamir: Phys.\ Rev.\ Lett.\ 71, 2691 (1993).
  \item V.~Furman and Y.~Shamir: Nucl.\ Phys.\ B 439, 54 (1995).
  \item P.~M.~Vranas: Phys.\ Rev.\ D 57, 1415 (1998).
  \item H.~Neuberger: Phys.\ Rev.\ D 57, 5417 (1998).
  \item Y.~Kikukawa and T.~Noguchi (unpublished) arXiv:hep-lat/9902022
        (1999).
  \item C.~R.~Allton et al.: Phys.\ Rev.\ D 78, 114509 (2008).
  \item R.~Frezzotti, P.~Grassi, S.~Sint, and P.~Weisz: Nucl.\ Phys.\ (Proc.\
        Suppl.) 83, 941 (2000).
  \item R.~Frezzotti, P.~Grassi, S.~Sint, and P.~Weisz: JHEP 0108, 058 (2001).
  \item R.~Frezzotti, S.~Sint, and P.~Weisz: JHEP 0107, 048 (2001).
  \item S.~Sint: in \emph{Perspectives in Lattice QCD}, Proceedings of the
        Workshop in Nara, Japan, 31 Oct--11 Nov 2005, edited by Y.~Kuramashi,
        p.~169 (World Scientific, Singapore 2007).
  \item A.~Shindler: Phys.\ Rep.\ 461, 37 (2008).
  \item R.~Frezzotti and G.~C.~Rossi: Nucl.\ Phys.\ (Proc.\ Suppl.) 128, 193
        (2004).
  \item C.~Pena, S.~Sint, and A.~Vladikas: JHEP 09, 069 (2004).
  \item C.~Gattringer and S.~Solbrig: Phys.\ Lett.\ B 621, 195 (2005).
  \item R.~Frezzotti and G.~C.~Rossi: JHEP 0408, 007 (2004).
  \item C.~Amsler et al.\ [Particle Data Group]: Phys.\ Lett.\ B 667, 1
        (2008).
  \item E.~Eichten and B.~Hill: Phys.\ Lett.\ B 234, 511 (1990).
  \item E.~Eichten and B.~Hill: Phys.\ Lett.\ B 240, 193 (1990).
  \item B.~A.~Thacker and G.~P.~Lepage: Phys.\ Rev.\ D 43, 196 (1991).
  \item G.~P.~Lepage et al.: Phys.\ Rev.\ D 46, 4052 (1992).
  \item M.~Neubert: Phys.\ Rept.\ 245, 259 (1994).
  \item A.~V.~Manohar and M.~B.~Wise: \emph{Heavy Quark Physics} (Cambridge
        University Press, Cambridge 2000).
  \item J.~D.~Bjorken and S.~D.~Drell: \emph{Relativistic Quantum Mechanics}\\
        (McGraw-Hill, New York 1964).
  \item J.~Heitger and R.~Sommer: JHEP 0402, 022 (2004).
  \item R.~Sommer: in \emph{Perspectives in Lattice QCD}, Proceedings of the
        Workshop in Nara, Japan, 31 Oct--11 Nov 2005, edited by Y.~Kuramashi,
        p.~209 (World Scientific, Singapore 2007).
  \item A.~Kronfeld: Nucl.\ Phys.\ B (Proc.\ Suppl.) 129, 46 (2004).
  \item T.~Onogi: PoS LAT2006, 017 (2006).
  \item A.~X.~El-Khadra, A.~S.~Kronfeld, and P.~B.~Mackenzie: Phys.\ Rev.\ D
        55, 3933 (1997).
  \item M.~L\"uscher, P.~Weisz, and U.~Wolff: Nucl.\ Phys.\ B 359, 221 (1991).
  \item M.~Della Morte, N.~Garron, M.~Papinutto, and R.~Sommer: JHEP 01, 007
        (2007).
  \item M.~Della Morte: PoS LATTICE2007, 008 (2007).
  \item J.~Heitger: Nucl.\ Phys.\ B (Proc.\ Suppl.) 181--182, 156 (2008).
  \item E.~Gamiz: PoS LATTICE2008, 014 (2009).
\end{enumerate}
